% !PRINTED 597
\weberpage{597}
%<<CONT
in $\Omega$ ist, so ist $Q$ die Gruppe derselben Gleichung in $\Omega(\varepsilon)$, oder
was dasselbe ist, in $\Omega(\varepsilon_1), \Omega(\varepsilon_2) \ldots \Omega(\varepsilon_{m-1})$. Nun können wir
aber den Schlusssatz in \S.\,157 anwenden. Nach diesem Satze muss
der Index $j$ des Theilers $Q$ von $P$ ein Theiler von $m$ sein, und
da $m$ als Primzahl vorausgesetzt ist, so ist $m = j$. Ausserdem ist
nach demselben Satze $\varepsilon$ rational durch die Wurzeln der Gleichung
$f(x) = 0$ darstellbar:

\[
\varepsilon = \psi(x_0, x_1 \ldots x_{m-1}).
\]

Diese Function gehört zur Gruppe $Q$, und wenn wir darauf
sämmtliche Permutationen der Gruppe $P$ anwenden, so erhalten
wir die Functionen $\varepsilon, \varepsilon_1, \varepsilon_2 \ldots \varepsilon_{m-1}$ und keine anderen. Diese
Functionen gehören zu den conjugirten Gruppen $\pi^{-1}Q\pi$. Da
aber jede dieser Functionen rational durch jede andere ausdrückbar ist, so müssen sie alle zu derselben Gruppe gehören,
d. h. $Q$ ist ein Normaltheiler von $P$.

Wir haben also hiermit den ersten Satz bewiesen:

\begin{quote}
\noindent I.\ \gesperrt{Wenn die Gruppe einer Gleichung $P$ durch
Adjunction der Wurzeln einer Abel'schen Gleichung reducirt wird, so hat $P$ einen Normaltheiler $Q$ von Primzahlindex.}
\end{quote}

Dieser Satz lässt sich auch umkehren.

Wenn nämlich die Gruppe $P$ einen Normaltheiler $Q$ vom
Index $m$ hat, so können wir eine zu $Q$ gehörige Function $\psi$
wählen, und die damit conjugirten Functionen $\psi, \psi_1, \psi_2 \ldots \psi_{m-1}$
gehören alle zu derselben Gruppe. Der Körper $\Omega(\psi)$ ist ein
Normalkörper und $\psi$ die Wurzel einer Normalgleichung. Im
Körper $\Omega(\psi)$ ist $Q$ die Gruppe von $f(x) = 0$ (\S.\,155). Wenn
aber $m$ eine Primzahl ist, so ist $\psi$ nach \S.\,163 die Wurzel einer
cyklischen Gleichung, und damit ist also bewiesen:

\begin{quote}
\noindent II.\ \gesperrt{Wenn die Gruppe $P$ von $f(x) = 0$ einen Normaltheiler $Q$ von Primzahlindex $m$ besitzt, so wird
durch Adjunction der Wurzel einer cyklischen
Gleichung $m^{\text{ten}}$ Grades die Gruppe $P$ auf $Q$
reducirt.}
\end{quote}

\begin{center}
\S.\,176.

\gesperrt{Metacyklische Gleichungen.}
\end{center}

Wir wollen eine Gleichung, deren vollständige Lösung sich
auf eine Kette von cyklischen Gleichungen zurückführen lässt,
%CONT>>
